% Converted from reference/DB4OTutorial (MadCap Flare, Java target)

\chapter{Introduction}
\label{tut:welcome}

Welcome to the db4o tutorial. In this tutorial we explore the basics of db4o and how to work with it.

\section*{Getting db4o}

db4o is maintained at \href{https://github.com/blakemcbride/db4o}{\nolinkurl{github.com/blakemcbride/db4o}}. Clone or download the repository and build it as described in its \href{https://github.com/blakemcbride/db4o/blob/master/BUILD.md}{\nolinkurl{BUILD.md}} file. The build needs JDK 17 or newer and nothing else:

\begin{lstlisting}
git clone https://github.com/blakemcbride/db4o.git
cd db4o
./bld jar
\end{lstlisting}

This writes the db4o jars to the target directory.

\section*{Adding db4o to a Project}

For this tutorial, add the all-in-one jar to your project:

\begin{itemize}
\item target/db4o-8.1-all.jar
\end{itemize}

A detailed overview of the jars is available in the reference documentation.

\section*{db4o Area of Use}

db4o is intended for embedded use with smaller databases around 2-16 GByte. As a general rule, if you expect your database to grow beyond 16 gigabytes, you should look at a larger database system.

db4o is explicitly single-threaded. Concurrent accesses will be synchronized against a global database lock. That means db4o cannot deal with a highly concurrent access, since it blocks on all operations. When you expect a high load and concurrent access, then you should consider a larger database system.

